%
% architecture.tex
%
% Ethos: Technical Architecture Specification
% Written for engineers integrating identity into AI agent systems.
%
% Compile: pdflatex architecture.tex
%

\documentclass[a4paper,10pt]{article}
\usepackage[margin=1in]{geometry}
\usepackage{booktabs}
\usepackage{listings}
\usepackage{xcolor}
\usepackage{enumitem}
\usepackage{longtable}
\usepackage[colorlinks=true,linkcolor=blue,citecolor=blue,urlcolor=blue]{hyperref}

\lstset{
  basicstyle=\ttfamily\small,
  keywordstyle=\color{blue!80!black},
  commentstyle=\color{green!50!black},
  stringstyle=\color{red!60!black},
  breaklines=true,
  frame=single,
  framerule=0.3pt,
  rulecolor=\color{gray!50},
  backgroundcolor=\color{gray!5},
  columns=fullflexible,
}

\newcommand{\pkg}[1]{\texttt{#1}}
\newcommand{\tool}[1]{\texttt{#1}}
\newcommand{\field}[1]{\texttt{#1}}
\newcommand{\cmd}[1]{\texttt{#1}}
\newcommand{\fpath}[1]{\texttt{#1}}

\begin{document}

\title{Ethos: Technical Architecture Specification}
\author{Identity Binding for Humans and AI Agents}
\date{March 2026 --- v0.3.3}
\hypersetup{
  pdftitle={Ethos: Technical Architecture Specification},
  pdfauthor={Punt Labs},
  pdfsubject={Architecture},
  pdfcreator={pdflatex}
}
\maketitle

\tableofcontents
\newpage

% ===================================================================
\section{Overview}
% ===================================================================

Ethos unifies a name, voice, email, GitHub handle, writing style,
personality, and skills into a single identity that other tools
optionally read.  It is written in Go and ships as a CLI binary plus
a Claude Code plugin with an MCP server.

The system serves two roles:

\begin{enumerate}
  \item \textbf{Identity registry} --- CRUD operations on identity
    YAML files stored at \fpath{\~{}/.punt-labs/ethos/identities/}.
    Humans and agents share the same schema.
  \item \textbf{Session roster} --- multi-participant awareness for
    Claude Code sessions.  Tracks who is present (human, primary
    agent, subagents), their personas, and parent-child relationships.
\end{enumerate}

Ethos is a \emph{sidecar}, not a dependency.  Other tools (Vox,
Beadle, Biff) read ethos state from known filesystem paths.  They do
not import ethos.  The file format is the contract.

\subsection{Design Principles}

\begin{enumerate}[label=\textbf{P\arabic*}.]
  \item \textbf{Sidecar, not dependency.}  Consumers read files at
    known paths.  No import, no version lock-step, no daemon required.
    Tools adopt ethos at their own pace.
  \item \textbf{Same schema for humans and agents.}  The \field{kind}
    field (\texttt{human} or \texttt{agent}) is the only structural
    difference.  An identity is an identity.
  \item \textbf{Agent definition is a channel binding.}  Like voice
    or email --- the \texttt{.md} file defines tools and workflow,
    ethos defines \emph{who}.
  \item \textbf{Generic extensions, not consumer fields.}  Ethos
    never adds fields for Vox, Beadle, or Biff.  Consumers own their
    own namespace files in the extension directory.
  \item \textbf{Fast.}  Go's $\sim$10\,ms cold start (vs Python's
    $\sim$200\,ms) keeps hook latency invisible.  Ethos is queried on
    session start, on every \cmd{/who}, and before every TTS call.
  \item \textbf{No daemon.}  All concurrency uses \texttt{flock}
    on lock files.  No background process, no database.
\end{enumerate}

% ===================================================================
\section{Package Map}
% ===================================================================

All reusable library code lives under \pkg{internal/}; the public
CLI entry point lives in \pkg{cmd/ethos}.  Module path:
\texttt{github.com/punt-labs/ethos}.

\begin{center}
\begin{longtable}{@{}lp{8cm}@{}}
\toprule
\textbf{Package} & \textbf{Responsibility} \\
\midrule
\endfirsthead
\toprule
\textbf{Package} & \textbf{Responsibility} \\
\midrule
\endhead
\pkg{cmd/ethos/}        & CLI entry point: subcommand dispatch, global
                          \cmd{--json} flag, interactive identity creation \\
\pkg{internal/identity/} & Core identity model (\texttt{Identity} struct),
                          validation, CRUD via \texttt{Store}, extension
                          I/O \\
\pkg{internal/session/}  & Session roster model (\texttt{Roster},
                          \texttt{Participant}), \texttt{flock}-based
                          concurrent store, PID-keyed session discovery,
                          staleness detection and purge \\
\pkg{internal/process/}  & Process tree walker: find topmost \texttt{claude}
                          ancestor PID via \texttt{ps}; result cached for
                          process lifetime \\
\pkg{internal/resolve/}  & Identity resolution chain: repo-local config
                          $\to$ global active identity $\to$ error \\
\pkg{internal/mcp/}      & MCP tool definitions and handlers (12~tools);
                          \texttt{Handler} struct with injected stores \\
\bottomrule
\end{longtable}
\end{center}

\subsection{Dependency Graph}

Dependencies flow strictly downward.  No package imports a package
above it in this ordering.

\begin{verbatim}
          cmd/ethos
        /   |   |   \       \
       /    |   |    \       \
    mcp  resolve process identity session
   /   \
identity session
\end{verbatim}

\pkg{cmd/ethos} imports all five internal packages directly.
\pkg{internal/mcp} depends on \pkg{identity} and \pkg{session}.
\pkg{internal/identity} and \pkg{internal/session} are independent ---
neither imports the other.  \pkg{internal/process} and
\pkg{internal/resolve} are leaf packages with no internal dependencies.

% ===================================================================
\section{Identity Model}
% ===================================================================

\subsection{Schema}

An identity is a YAML file at
\fpath{\~{}/.punt-labs/ethos/identities/<handle>.yaml}.

\begin{center}
\begin{longtable}{@{}llp{6cm}@{}}
\toprule
\textbf{Field} & \textbf{Type} & \textbf{Description} \\
\midrule
\endfirsthead
\toprule
\textbf{Field} & \textbf{Type} & \textbf{Description} \\
\midrule
\endhead
\field{name}          & string (required)  & Display name \\
\field{handle}        & string (required)  & Unique lookup key, used as filename \\
\field{kind}          & string (required)  & \texttt{human} or \texttt{agent} \\
\field{email}         & string             & Beadle channel binding \\
\field{github}        & string             & Biff channel binding \\
\field{voice}         & object             & Vox channel binding:
                                             \field{provider} and \field{voice\_id} \\
\field{agent}         & string             & Path to Claude Code agent \texttt{.md}
                                             file \\
\field{writing\_style} & multiline string  & Prose style directives \\
\field{personality}   & multiline string   & Behavioral directives \\
\field{skills}        & list of strings    & Capability tags \\
\bottomrule
\end{longtable}
\end{center}

\subsection{Validation}

\texttt{Identity.Validate()} enforces:

\begin{itemize}
  \item \field{name} and \field{handle} are non-empty.
  \item \field{handle} matches
    \verb|^[a-z0-9]([a-z0-9-]*[a-z0-9])?$| (lowercase alphanumeric,
    may contain internal hyphens but not leading/trailing, used as
    filename).
  \item \field{kind} is one of \texttt{human} or \texttt{agent}.
  \item If \field{voice\_id} is set, \field{voice.provider} is
    required.
\end{itemize}

Validation errors are returned as a structured
\texttt{ValidationError} type, not strings.

\subsection{Persistence}

\texttt{Store} takes a root directory path
(\fpath{\~{}/.punt-labs/ethos} by default) and provides:

\begin{center}
\begin{tabular}{@{}lp{7cm}@{}}
\toprule
\textbf{Method} & \textbf{Behavior} \\
\midrule
\texttt{Save(id)}       & Atomic create with \texttt{O\_EXCL} --- rejects
                          duplicate handles.  Creates the extension
                          directory. \\
\texttt{Load(handle)}   & Read YAML, assemble extensions from
                          \fpath{<handle>.ext/} directory. \\
\texttt{List()}          & Return all identities (calls \texttt{Load}
                          per handle, includes extensions). \\
\texttt{Active()}        & Read handle from the \fpath{active} file,
                          then \texttt{Load()} the full identity. \\
\texttt{SetActive(h)}   & Write the handle to the \fpath{active} file. \\
\texttt{Exists(handle)}  & Check file existence. \\
\bottomrule
\end{tabular}
\end{center}

% ===================================================================
\section{Extension System}
% ===================================================================

Consumers of ethos (Beadle, Biff, Vox) store tool-specific attributes
on identities without modifying the ethos schema.  Extensions use a
namespaced, file-per-tool design.

\subsection{Storage Layout}

\begin{verbatim}
~/.punt-labs/ethos/identities/
  mal.yaml                     # ethos owns -- core identity fields
  mal.ext/                     # extension directory
    beadle.yaml                # beadle owns -- GPG key, IMAP config
    biff.yaml                  # biff owns -- preferred TTY
    vox.yaml                   # vox owns -- default mood
\end{verbatim}

Each file is a flat YAML map of key-value pairs.  Ethos never reads
or interprets extension contents except to assemble the merged view
returned by \texttt{Load()} and the \tool{get\_identity} MCP tool.

\subsection{Ownership Rules}

\begin{enumerate}
  \item Ethos creates the \fpath{<handle>.ext/} directory on identity
    creation.
  \item Each tool manages its own \fpath{<namespace>.yaml} file.
  \item Tools may read/write their file directly (sidecar contract) or
    through the ethos CLI/MCP interface.
  \item Tools must not write to another tool's namespace file.
\end{enumerate}

\subsection{Validation Constraints}

\begin{center}
\begin{tabular}{@{}llr@{}}
\toprule
\textbf{Field} & \textbf{Pattern} & \textbf{Limit} \\
\midrule
Namespace   & \verb|^[a-z][a-z0-9-]*$|    & 32 chars \\
Key         & \verb|^[a-z][a-z0-9_]*$|    & 64 chars \\
Value       & Any YAML scalar             & 4096 bytes \\
Keys per namespace    & ---               & 64 \\
Namespaces per persona & ---              & 32 \\
\bottomrule
\end{tabular}
\end{center}

\subsection{Two Scopes}

Extensions exist at two independent scopes with different types:

\begin{enumerate}
  \item \textbf{Persona-level} (durable) --- files in
    \fpath{<handle>.ext/}.  Flat \texttt{map[string]string} (YAML
    scalars only).  Persist across sessions.
  \item \textbf{Session-participant-level} (ephemeral) --- the
    \field{ext} field on each \texttt{Participant} in the session
    roster.  Typed as \texttt{map[string]any} (arbitrary nested
    values).  Deleted when the session ends.
\end{enumerate}

The two scopes are structurally different.  Persona-level extensions
enforce flat string values; session-level extensions accept arbitrary
YAML structures.  A tool may store durable defaults at the persona
level and ephemeral session state at the participant level.

% ===================================================================
\section{Session Roster}
% ===================================================================

The session roster provides multi-participant awareness for Claude
Code sessions.  It answers two questions: ``Who am I?'' and ``Who is
everyone else?''

\subsection{Data Model}

Stored at \fpath{\~{}/.punt-labs/ethos/sessions/<session-id>.yaml}:

\begin{center}
\begin{longtable}{@{}llp{5.5cm}@{}}
\toprule
\textbf{Field} & \textbf{Source} & \textbf{Purpose} \\
\midrule
\endfirsthead
\toprule
\textbf{Field} & \textbf{Source} & \textbf{Purpose} \\
\midrule
\endhead
\field{session}     & Hook input        & Session identifier \\
\field{started}     & Roster creation   & RFC~3339 timestamp \\
\field{participants} & Hook lifecycle   & Ordered list of \texttt{Participant}
                                          records \\
\bottomrule
\end{longtable}
\end{center}

Each participant record:

\begin{center}
\begin{longtable}{@{}llp{5cm}@{}}
\toprule
\textbf{Field} & \textbf{Type} & \textbf{Description} \\
\midrule
\endfirsthead
\toprule
\textbf{Field} & \textbf{Type} & \textbf{Description} \\
\midrule
\endhead
\field{agent\_id}    & string  & OS login, Claude PID, or subagent ID \\
\field{persona}      & string  & Ethos handle; null when no identity \\
\field{agent\_type}  & string  & Raw type (e.g., \texttt{Explore},
                                 \texttt{code-reviewer}) \\
\field{parent}       & string  & Parent's \field{agent\_id}; null for root \\
\field{ext}          & \texttt{map[string]any}  & Tool-scoped ephemeral metadata \\
\bottomrule
\end{longtable}
\end{center}

\subsection{Tree Structure}

Participants form a tree.  The root (no parent) is the human user.
The primary Claude agent is the root's child.  Subagents are children
of the primary agent.

\begin{verbatim}
mal (root -- $USER)
  +-- 19147 (archie -- Claude PID)
       +-- a5734dd (code-reviewer)
       +-- a93c2be (silent-failure-hunter)
       +-- b8823ff (Explore, no persona)
\end{verbatim}

Any participant can derive relationships from the tree: root (walk up
to \field{parent:~null}), initiator (my \field{parent}), delegates
(anyone whose \field{parent} is my \field{agent\_id}), siblings (same
\field{parent}).

\subsection{Lifecycle via Hooks}

\begin{center}
\begin{tabular}{@{}llp{5cm}@{}}
\toprule
\textbf{Hook Event} & \textbf{Roster Operation} & \textbf{Data Available} \\
\midrule
\texttt{SessionStart}  & Create roster, join root + primary
                        & \field{session\_id}, \texttt{\$USER},
                          \texttt{\$PPID} as Claude PID \\
\texttt{SubagentStart}  & Join subagent
                        & \field{session\_id}, \field{agent\_id},
                          \field{agent\_type} \\
\texttt{SubagentStop}   & Leave subagent
                        & \field{session\_id}, \field{agent\_id} \\
\texttt{SessionEnd}     & Delete roster + current session file
                        & \field{session\_id} \\
\bottomrule
\end{tabular}
\end{center}

\subsection{Concurrency}

Multiple \texttt{SubagentStart} hooks may fire in parallel.  Roster
writes use \texttt{flock(LOCK\_EX)} on a per-session lock file
(\fpath{sessions/<session-id>.lock}) before read-modify-write.  Roster
files are written atomically via temp file + rename.

\subsection{Session Discovery}

Non-hook callers (Biff, Vox) need the session ID without receiving it
on stdin.  The \texttt{SessionStart} hook writes the session ID to a
PID-keyed file:

\begin{verbatim}
~/.punt-labs/ethos/sessions/current/<claude-pid>
\end{verbatim}

Any descendant process walks the process tree to the topmost
\texttt{claude} ancestor PID, reads that file, and gets the session
ID.  The \texttt{SessionEnd} hook deletes the file.

\subsection{Staleness Detection}

\texttt{Store.Purge()} cleans up rosters from crashed sessions:

\begin{enumerate}
  \item Extract the primary agent's \field{agent\_id} (second
    participant).
  \item If numeric: check whether the PID is alive via
    \texttt{os.FindProcess} + signal~0.  If dead, treat as stale
    immediately (no TTL fallback).
  \item If non-numeric (parse fails): age-based fallback (24-hour
    TTL on the \field{started} timestamp).
  \item Rosters with fewer than 2~participants are always stale
    (malformed).
\end{enumerate}

% ===================================================================
\section{Identity Resolution}
% ===================================================================

\texttt{resolve.Resolve(repoRoot)} determines the active identity
handle through a two-level priority chain:

\begin{enumerate}
  \item \textbf{Repo-local config} ---
    \fpath{.punt-labs/ethos/config.yaml} in the repository root.
    Contains a \texttt{RepoConfig} struct with \field{active} and
    \field{agents} fields.  If \field{active} is non-empty, return it.
  \item \textbf{Global active identity} ---
    \fpath{\~{}/.punt-labs/ethos/active} (plain text file).  If
    non-empty, return it.
  \item \textbf{Error} --- no active identity configured.
\end{enumerate}

This mirrors Git's resolution: \fpath{.git/config} overrides
\fpath{\~{}/.gitconfig}.

% ===================================================================
\section{Process Tree Walker}
% ===================================================================

\texttt{process.FindClaudePID()} identifies the Claude Code instance
that spawned the current process.  It is used to determine the
primary agent's \field{agent\_id} for session rosters and to locate
PID-keyed session files.

\begin{enumerate}
  \item Run \texttt{ps -eo pid=,ppid=,comm=} to snapshot the process
    table.
  \item Parse into a map of PID $\to$ (parent PID, command name).
  \item Walk from the current PID upward, following parent pointers.
  \item Track the topmost PID where the command name is
    \texttt{claude}.
  \item Cap the walk at 10~iterations to prevent infinite loops in
    degenerate process trees.
  \item If no \texttt{claude} ancestor is found, fall back to
    \texttt{os.Getppid()}.
  \item Cache the result for the lifetime of the process (the answer
    does not change).
\end{enumerate}

% ===================================================================
\section{Storage Layout}
% ===================================================================

\begin{center}
\begin{longtable}{@{}lll@{}}
\toprule
\textbf{Scope} & \textbf{Path} & \textbf{Git-tracked?} \\
\midrule
\endfirsthead
\toprule
\textbf{Scope} & \textbf{Path} & \textbf{Git-tracked?} \\
\midrule
\endhead
Identities      & \fpath{\~{}/.punt-labs/ethos/identities/<handle>.yaml}
                 & No \\
Extensions       & \fpath{\~{}/.punt-labs/ethos/identities/<handle>.ext/<ns>.yaml}
                 & No \\
Active identity  & \fpath{\~{}/.punt-labs/ethos/active}
                 & No \\
Sessions         & \fpath{\~{}/.punt-labs/ethos/sessions/<id>.yaml}
                 & No \\
Session locks    & \fpath{\~{}/.punt-labs/ethos/sessions/<id>.lock}
                 & No \\
Current session  & \fpath{\~{}/.punt-labs/ethos/sessions/current/<pid>}
                 & No \\
Repo config      & \fpath{.punt-labs/ethos/config.yaml}
                 & Yes \\
Repo agents      & \fpath{.punt-labs/ethos/agents/<name>.yaml}
                 & Yes \\
\bottomrule
\end{longtable}
\end{center}

All user-global state lives under \fpath{\~{}/.punt-labs/ethos/}.
Repo-scoped state lives under \fpath{.punt-labs/ethos/} in the
repository root.

% ===================================================================
\section{CLI Surface}
% ===================================================================

The binary exposes subcommands grouped by concern.  All commands
accept a global \cmd{--json} flag for machine-readable output.

\subsection{Identity Commands}

\begin{center}
\begin{tabular}{@{}lp{7cm}@{}}
\toprule
\textbf{Command} & \textbf{Description} \\
\midrule
\cmd{whoami [handle]}     & Show or set the active identity. \\
\cmd{create [-f FILE]}    & Create identity interactively or from a
                            YAML file. \\
\cmd{list}                & List all identities; active marked with
                            \texttt{*}. \\
\cmd{show <handle>}       & Show full identity details including
                            extensions. \\
\bottomrule
\end{tabular}
\end{center}

\subsection{Extension Commands}

\begin{center}
\begin{tabular}{@{}lp{7cm}@{}}
\toprule
\textbf{Command} & \textbf{Description} \\
\midrule
\cmd{ext get <p> <ns> [key]}       & Read one key or all keys in a
                                     namespace. \\
\cmd{ext set <p> <ns> <key> <val>} & Write a key-value pair. \\
\cmd{ext del <p> <ns> [key]}       & Delete one key or an entire
                                     namespace. \\
\cmd{ext list <p>}                 & List all namespaces for a
                                     persona. \\
\bottomrule
\end{tabular}
\end{center}

\subsection{Session Commands}

\begin{center}
\begin{tabular}{@{}lp{7cm}@{}}
\toprule
\textbf{Command} & \textbf{Description} \\
\midrule
\cmd{iam <persona>}            & Declare persona in current session
                                 (auto-detects session ID and agent
                                 ID). \\
\cmd{session}                  & Show current session roster. \\
\cmd{session create}           & Create a new roster (called by
                                 SessionStart hook). \\
\cmd{session join}             & Add a participant (called by
                                 SubagentStart hook). \\
\cmd{session leave}            & Remove a participant (called by
                                 SubagentStop hook). \\
\cmd{session write-current}    & Write PID-keyed session file. \\
\cmd{session delete-current}   & Delete PID-keyed session file. \\
\cmd{session purge}            & Clean up stale rosters. \\
\bottomrule
\end{tabular}
\end{center}

\subsection{Admin Commands}

\begin{center}
\begin{tabular}{@{}lp{7cm}@{}}
\toprule
\textbf{Command} & \textbf{Description} \\
\midrule
\cmd{version}              & Print version string. \\
\cmd{doctor}               & Check installation health (identity
                             directory, active identity). \\
\cmd{serve}                & Start MCP server (stdio transport). \\
\cmd{uninstall [--purge]}  & Remove plugin; \cmd{--purge} also
                             removes binary and data. \\
\bottomrule
\end{tabular}
\end{center}

% ===================================================================
\section{MCP Tool Surface}
% ===================================================================

\texttt{ethos serve} starts a stdio-transport MCP server that
registers 12~tools via \texttt{Handler.RegisterTools()}.  The
\texttt{Handler} struct receives injected \texttt{identity.Store}
and \texttt{session.Store} instances.

\subsection{Identity Tools}

\begin{center}
\begin{longtable}{@{}llp{5.5cm}@{}}
\toprule
\textbf{Tool} & \textbf{Required Params} & \textbf{Description} \\
\midrule
\endfirsthead
\toprule
\textbf{Tool} & \textbf{Required Params} & \textbf{Description} \\
\midrule
\endhead
\tool{whoami}           & ---                 & Show active identity; optional
                                                \field{handle} to set it. \\
\tool{list\_identities} & ---                 & List all identities with active
                                                status. \\
\tool{get\_identity}    & \field{handle}       & Return full identity including
                                                extensions. \\
\tool{create\_identity} & \field{name},
                          \field{handle},
                          \field{kind}         & Create identity from provided
                                                fields. \\
\bottomrule
\end{longtable}
\end{center}

\subsection{Extension Tools}

\begin{center}
\begin{longtable}{@{}llp{5.5cm}@{}}
\toprule
\textbf{Tool} & \textbf{Required Params} & \textbf{Description} \\
\midrule
\endfirsthead
\toprule
\textbf{Tool} & \textbf{Required Params} & \textbf{Description} \\
\midrule
\endhead
\tool{ext\_get}  & \field{persona}, \field{namespace}
                 & Read extension keys (optional \field{key} for
                   single lookup). \\
\tool{ext\_set}  & \field{persona}, \field{namespace},
                   \field{key}, \field{value}
                 & Write a key-value pair. \\
\tool{ext\_del}  & \field{persona}, \field{namespace}
                 & Delete a key or entire namespace (optional
                   \field{key}). \\
\tool{ext\_list} & \field{persona}
                 & List all extension namespaces. \\
\bottomrule
\end{longtable}
\end{center}

\subsection{Session Tools}

\begin{center}
\begin{longtable}{@{}llp{5.5cm}@{}}
\toprule
\textbf{Tool} & \textbf{Required Params} & \textbf{Description} \\
\midrule
\endfirsthead
\toprule
\textbf{Tool} & \textbf{Required Params} & \textbf{Description} \\
\midrule
\endhead
\tool{session\_iam}     & \field{session\_id},
                          \field{agent\_id},
                          \field{persona}
                        & Declare persona for a participant. \\
\tool{session\_roster}  & \field{session\_id}
                        & Return full roster with tree. \\
\tool{session\_join}    & \field{session\_id},
                          \field{agent\_id}
                        & Register a participant (optional
                          \field{persona}, \field{parent},
                          \field{agent\_type}). \\
\tool{session\_leave}   & \field{session\_id},
                          \field{agent\_id}
                        & Deregister a participant. \\
\bottomrule
\end{longtable}
\end{center}

Most tools return JSON objects on success.  Mutating tools
(\tool{whoami} when setting, \tool{ext\_set}, \tool{ext\_del},
\tool{session\_iam}, \tool{session\_leave}) return plain-text
confirmation strings instead.  Error responses are also plain strings.

% ===================================================================
\section{Design Invariants}
% ===================================================================

These invariants are enforced by the architecture and hold across
all operations:

\begin{enumerate}
  \item \textbf{Sidecar contract.}  No consumer imports ethos.
    The contract is file paths and YAML formats.  Upgrading ethos
    never forces consumer upgrades.

  \item \textbf{Unified schema.}  Humans and agents use the same
    \texttt{Identity} struct.  The \field{kind} field is the only
    difference.  No separate code paths.

  \item \textbf{No consumer-specific fields.}  Ethos will never add
    \field{gpg\_key\_id}, \field{preferred\_tty}, or any field that
    serves a single consumer.  The extension system exists for this.

  \item \textbf{Atomic identity creation.}  \texttt{Store.Save}
    opens the file with \texttt{O\_EXCL}.  Two concurrent creates
    of the same handle: one succeeds, one fails.  No partial writes.

  \item \textbf{Atomic roster writes.}  Roster files are written to
    a temp file, then renamed.  No reader sees a partial roster.

  \item \textbf{Flock concurrency.}  All roster mutations acquire
    \texttt{flock(LOCK\_EX)} on a per-session lock file.  No daemon,
    no database, no race conditions on concurrent subagent hooks.

  \item \textbf{Extension isolation.}  Each consumer's extension
    data lives in its own YAML file.  Concurrent writes by different
    tools cannot corrupt each other's data.

  \item \textbf{Process tree caching.}  \texttt{FindClaudePID()}
    caches the result for the process lifetime.  The Claude ancestor
    PID does not change during a session.

  \item \textbf{No shell execution.}  All subprocess invocations use
    argument lists, never shell interpretation.

  \item \textbf{Depth-bounded tree walk.}  The process tree walk
    caps at 10~iterations to prevent infinite loops in degenerate
    process tables.
\end{enumerate}

% ===================================================================
\section{Security Boundaries}
% ===================================================================

\subsection{Filesystem Boundary}

\begin{itemize}
  \item \textbf{User-global state}: \fpath{\~{}/.punt-labs/ethos/} ---
    identities, extensions, sessions.  Standard user-mode file
    permissions.
  \item \textbf{Repo-local state}: \fpath{.punt-labs/ethos/} ---
    config, agent roster.  Git-tracked; visible to collaborators.
  \item \textbf{No secrets.}  Ethos stores no passwords, API keys,
    or tokens.  Sensitive bindings (GPG keys, IMAP passwords) live
    in consumer extensions or credential stores, not in ethos.
\end{itemize}

\subsection{Process Boundary}

\begin{itemize}
  \item \textbf{MCP}: stdio transport.  No network listener.
  \item \textbf{ps}: read-only process table query for tree walking.
    Signal~0 for process liveness checks (no side effects).
\end{itemize}

\subsection{Trust Boundary}

\begin{itemize}
  \item \textbf{Identity data is self-declared.}  Ethos does not
    verify that the email, GitHub handle, or voice binding is
    authentic.  Verification is the responsibility of consumers
    (Beadle verifies PGP, Biff verifies GitHub tokens).
  \item \textbf{Extension data is opaque.}  Ethos enforces structural
    constraints (namespace format, key format, size limits) but does
    not validate or interpret values.
\end{itemize}

% ===================================================================
\section{Key Properties}
% ===================================================================

The architecture guarantees these properties, verifiable by
inspection and testing:

\begin{enumerate}
  \item \textbf{Resolution totality}: \texttt{Resolve()} always
    returns a handle or an error.  There is no silent empty-string
    fallback.

  \item \textbf{Extension namespace isolation}: concurrent writes by
    different tools target different files.  No merge conflicts.

  \item \textbf{Session liveness}: \texttt{Purge()} uses process
    signal~0 for numeric PIDs and age-based TTL for non-numeric IDs.
    Every stale roster is eventually cleaned up.

  \item \textbf{Hook-driven lifecycle}: the session roster is fully
    managed by four Claude Code hook events.  No polling, no daemon,
    no manual intervention.

  \item \textbf{Cold start performance}: Go binary with no runtime
    dependencies.  $\sim$10\,ms cold start measured on Apple~M2.
    Hooks are invisible to the user.

  \item \textbf{Store injectability}: all CRUD operations take a
    \texttt{Store} with a configurable root path.  Tests use
    \texttt{t.TempDir()}.  No hardcoded \texttt{\$HOME}.
\end{enumerate}

\end{document}
